\section{Benchmark eligibility}
\label{sec:eligibility}

\subsection{Where a numerical-floor criterion fails}

Our initial eligibility condition computed the ceiling from the numerical floor,
\begin{equation}
  \text{ceiling} = \log L_0 - \log\!\left(L_0 \cdot 100\varepsilon\right)
  \approx 31.44 \text{ nat} .
\end{equation}

By that criterion Rosenbrock ($d = 5$, standard starting point) was eligible: the
median log improvement of four baselines was $1.8175$~nat, leaving an apparent
margin of $29.62$~nat to the ceiling.

In fact all four baselines returned \emph{exactly} $1.8175$~nat. We identified the
cause: the basin of the standard starting point contains a strict local minimum.

The literature distinguishes \emph{two variants} of the extended Rosenbrock
function \citep{kok2009rosenbrock}. The uncoupled sum of two-dimensional problems,
defined only in even dimensions, has only the global minimum, whereas the
coordinate-coupled variant acquires additional stationary points for $d \geq 4$
\citep{shang2006rosenbrock,kok2009rosenbrock}. We use the latter:
\begin{equation}
  L(x) = \sum_{i=1}^{d-1}
    \left[ 100\,(x_{i+1} - x_i^2)^2 + (1 - x_i)^2 \right].
  \label{eq:rosenbrock}
\end{equation}

The point below was obtained by converging our own implementation at $d = 5$; it
is not transcribed from a published table. Its first coordinate,
$x_1 = -0.962$, is consistent with the approximate location of the local minimum
reported in the literature.

\begin{center}
\begin{tabular}{ll}
\toprule
Quantity & Value \\
\midrule
$x^\ast$ & $(-0.96205102,\ 0.93573939,\ 0.88071360,\ 0.77787767,\ 0.60509367)$ \\
Loss & $3.930839434133$ \\
$\lVert \nabla L \rVert$ & $1.06 \times 10^{-8}$ \\
Hessian & smallest eigenvalue $+0.595$ (positive definite) \\
\bottomrule
\end{tabular}
\end{center}

Converging L-BFGS from the standard starting point reaches this point rather than
the global minimum. The attainable log improvement is
$\log(24.2 / 3.930839) = 1.8175$~nat, exactly the observed baseline value. The
margin was therefore \emph{zero}.

\subsection{Correcting to an attainable ceiling}

\begin{align}
  L_{\text{ref}} &= \min \text{ over reference runs started from the task's own
    initial point}, \\
  J_{\text{achievable}} &= \log L_0 - \log \max(L_{\text{ref}}, L_{\text{floor}}).
\end{align}

We use a panel of reference solvers and forbid relying on a single one. Evidence
came from the $\kappa = 10^6$ quadratic:

\begin{center}
\begin{tabular}{lrrl}
\toprule
Solver & Final loss & $\lVert \nabla L \rVert$ & Status \\
\midrule
\texttt{lbfgs}  & $1.007485 \times 10^{-6}$  & $3.456 \times 10^{-2}$  & not converged \\
\texttt{newton} & $3.059391 \times 10^{-32}$ & $2.650 \times 10^{-16}$ & converged \\
\texttt{sgd}    & \texttt{nan}               & ---                     & failed \\
\bottomrule
\end{tabular}
\end{center}

Had we used L-BFGS alone, $J_{\text{achievable}}$ would have been underestimated
at roughly $28.3$~nat.

\emph{Results from a different initialization do not raise the ceiling.} The
controller always starts from the task's own initial point, so an optimum in
another basin is not attainable. We record such results as diagnostics only.
Starting from $(0.9, \ldots, 0.9)$, Rosenbrock at $d = 5$ reaches the global
minimum; feeding that value into $L_{\text{ref}}$ inflates
$J_{\text{achievable}}$ from $2.58$ to $31.44$, admitting a spec that is in fact
trapped in a local minimum.

Planner-family controllers are excluded from the panel. Including them would leak
planner results into spec selection.

\subsection{Detecting duplicated seeds}

Because \texttt{initial\_loss} values can coincide by chance, we also compare the
starting-point vectors. This check revealed that the default configuration of
Rosenbrock $d = 5$ did not randomize the starting point, so seeds 2, 3, and 4
produced \emph{the same instance}. Twenty-four runs were bitwise identical in
every column, and that diagnostic's nominal $n = 3$ was in fact $n = 1$.

\subsection{Final eligibility conditions and verdicts}

\begin{center}
\begin{tabular}{l}
\toprule
Condition \\
\midrule
\texttt{failure\_rate} $= 0$ \\
Joint floor-hit rate $\leq 1/3$ \\
Median $\JE \geq 1$~nat for every baseline \\
$J_{\text{achievable}} - \text{median } \JE \geq 3$~nat \\
Spread among converged reference solvers $\leq 0.5$~nat \\
Each seed yields a genuinely different instance \\
\bottomrule
\end{tabular}
\end{center}

\begin{table}[!ht]
  \centering
  \begin{tabular}{llrlr}
  \toprule
  Spec & Verdict & $J_{\text{achievable}}$ & Limiting factor & Minimum margin \\
  \midrule
  quad $d{=}100$, $\kappa{=}10^3$ & accepted & 31.44 & numerical floor & 13.71 \\
  quad $d{=}100$, $\kappa{=}10^4$ & accepted & 31.44 & numerical floor & 20.98 \\
  quad $d{=}100$, $\kappa{=}10^5$ & accepted & 31.44 & numerical floor & 21.88 \\
  quad $d{=}100$, $\kappa{=}10^6$ & accepted & 31.44 & numerical floor & 21.93 \\
  rosen $d{=}5$                   & rejected & 1.82  & critical point  & $-0.00$ \\
  rosen $d{=}5$ (random start)    & rejected & 2.58  & critical point  & 0.00 \\
  micro-neural full-batch         & accepted & 31.44 & numerical floor & 11.68 \\
  micro-neural minibatch          & accepted & 31.44 & numerical floor & 28.75 \\
  \bottomrule
  \end{tabular}
  \caption{Eligibility verdicts. The margin for the Rosenbrock family is exactly
  zero, so those specs cannot discriminate between controllers. We do not
  interpret this as a failure of the method on nonlinear problems.}
  \label{tab:eligibility}
\end{table}
